\subsection{Definition of the integral of a vector-valued function}

Recall that a mapping \(\mathbf{f}\) of X into E is said to be weakly continuous if, for every \(\mathbf{z}' \in \mathrm{E}'\) , the mapping \(x \mapsto \langle \mathbf{f}(x), \mathbf{z}' \rangle\) of X into \(\mathbf{C}\) (in other words the mapping \(\mathbf{z}' \circ \mathbf{f}\) , also denoted \(\langle \mathbf{f}, \mathbf{z}' \rangle\) ) is continuous. We shall say that a mapping \(\mathbf{f}\) of X into E is scalarly of compact support if, for every \(\mathbf{z}' \in \mathrm{E}'\) , the mapping \(x \mapsto \langle \mathbf{f}(x), \mathbf{z}' \rangle\) has compact support. We denote by \(\widetilde{\mathcal{K}}(\mathrm{X};\mathrm{E})\) the space of mappings of X into E that are weakly continuous and scalarly of compact support; it is clear that \(\widetilde{\mathcal{K}}(\mathrm{X};\mathrm{E}) \supset \mathcal{K}(\mathrm{X};\mathrm{E})\) , but these two spaces are not necessarily identical (see Example 2 below); they are, however, equal when E is finite-dimensional.

Note that in the definition of a function that is weakly continuous (resp. scalarly of compact support), the topology of E intervenes only through the intermediation of the dual \(E'\) of E; thus, the set of these functions is not changed when the topology of E is replaced by any locally convex topology for which the dual is the same.

If E and F are two vector spaces in duality, we note that it means the same to say that a mapping of X into E is continuous for \(\sigma(\mathrm{E},\mathrm{F})\) and to say that it is weakly continuous.

Let \(\mathbf{f}\) be a mapping of \(X\) into \(E\) , weakly continuous and scalarly of compact support, and let \(\mu\) be a measure on \(X\) ; for every \(\mathbf{z}' \in E'\) we have \(\mathbf{z}' \circ \mathbf{f} \in \mathcal{K}(X)\) ; set

\[
\varphi (\mathbf {z} ^ {\prime}) = \int \langle \mathbf {f} (x), \mathbf {z} ^ {\prime} \rangle d \mu (x) = \mu (\mathbf {z} ^ {\prime} \circ \mathbf {f}).
\]

It is clear that \(\varphi\) is a linear form on \(\mathbf{E}'\) , hence is an element of \(\mathbf{E}^{\prime*}\) .

DEFINITION 1. --- For every function \(\mathbf{f} \in \widetilde{\mathcal{K}}(\mathrm{X};\mathrm{E})\) we call integral of \(\mathbf{f}\) with respect to \(\mu\) , and denote by \(\int \mathbf{f} d\mu\) or \(\int \mathbf{f}(x) d\mu(x)\) , or \(\int \mathbf{f}\mu\) , or \(\int \mathbf{f}(x)\mu(x)\) , the element of \(\mathrm{E}'^*\) defined by

\[
\left\langle \int \mathbf {f} d \mu , \mathbf {z} ^ {\prime} \right\rangle = \int \left\langle \mathbf {f}, \mathbf {z} ^ {\prime} \right\rangle d \mu \quad \text {for all} \mathbf {z} ^ {\prime} \in \mathrm{E} ^ {\prime}.\tag{1}
\]

We note that even if E is Hausdorff and \(\mathbf{f} \in \mathcal{K}(\mathrm{X};\mathrm{E})\) , one does not necessarily have \(\int f d\mu \in E\) (Exer. 1; cf.~No.~3).

Examples. --- 1) Suppose that E is finite-dimensional over C and Hausdorff, so that if \((\mathbf{e}_{i})_{1\leqslant i\leqslant n}\) is a basis of E then the mapping

\[
\left(\xi_ {1}, \dots \xi_ {n}\right) \mapsto \sum_ {i = 1} ^ {n} \xi_ {i} \mathbf {e} _ {i}
\]

is an isomorphism of \(\mathbf{C}^n\) onto \(\mathbf{E}\) . We then know that every linear form on \(\mathbf{E}\) is continuous, in other words that \(\mathbf{E}'\) is identical to the algebraic dual \(\mathbf{E}^*\) of \(\mathbf{E}\) , and \(\mathbf{E}'^*\) may be canonically identified with \(\mathbf{E}\) . Let \((\mathbf{e}_i')_{1 \leqslant i \leqslant n}\) be the basis of \(\mathbf{E}'\) dual to \((\mathbf{e}_i)\) ; for a mapping \(\mathbf{f}\) of \(\mathbf{X}\) into \(\mathbf{E}\) to be weakly continuous and scalarly of compact support, it is necessary and sufficient that the functions \(f_i = \mathbf{e}_i' \circ \mathbf{f}\) be continuous with compact support; we then have \(\mathbf{f}(x) = \sum_{i=1}^{n} f_i(x) \mathbf{e}_i\) for all \(x \in \mathbf{X}\) , and

\[
\int \mathbf {f} d \mu = \sum_ {i = 1} ^ {n} \mu (f _ {i}) \mathbf {e} _ {i}.
\]

\begin{enumerate}
\def\labelenumi{\arabic{enumi})}
\setcounter{enumi}{1}
\tightlist
\item
  Let us take for E the space \(\mathcal{M}(\mathrm{X};\mathbf{C})\) of measures on X, equipped with the vague topology ( \(\S1\) , No.~9); the dual \(E'\) of E may then be canonically identified with the space \(\mathcal{K}(\mathrm{X};\mathbf{C})\) (TVS, II, \(\S6\) , No.~2, Prop. 3). The mapping \(x \mapsto \varepsilon_{x}\) of X into E is continuous ( \(\S1\) , No.~9, Prop. 13), but its support is not compact if X is not compact; however, it is scalarly of compact support, because for every function \(f \in E'\) the function \(x \mapsto \langle \varepsilon_x, f \rangle = f(x)\) by definition has compact support. Moreover,
\end{enumerate}

\[
\int \left\langle \varepsilon_ {x}, f \right\rangle d \mu = \int f (x) d \mu (x) = \left\langle \mu , f \right\rangle
\]

for every function \(f \in \mathcal{K}(\mathrm{X};\mathbf{C}) = \mathrm{E}'\) , which proves that

\[
\int \varepsilon_ {x} d \mu (x) = \mu
\]

for every measure \(\mu\) on X.

\begin{enumerate}
\def\labelenumi{\arabic{enumi})}
\setcounter{enumi}{2}
\tightlist
\item
  If \(\mathbf{E}\) is Hausdorff then, for every point \(y \in \mathbf{X}\) and every function \(\mathbf{f} \in \widetilde{\mathcal{K}}(\mathbf{X};\mathbf{E})\) , we have
\end{enumerate}

\[
\int \mathbf {f} d \varepsilon_ {y} = \mathbf {f} (y)
\]

because \(\int \langle \mathbf{f},\mathbf{z}^{\prime}\rangle d\varepsilon_{y} = \langle \mathbf{f}(y),\mathbf{z}^{\prime}\rangle\) by definition.

Remarks. --- 1) If E is a locally convex space and N is the closure of \(\{0\}\) in E, so that \(E_{1} = E/N\) is the Hausdorff locally convex space associated with E, we know that the duals \(E'\) and \(E_{1}'\) are identical; for a function f to belong to \(\widetilde{\mathcal{K}}(X;E)\) , it is necessary and sufficient that \(f_{1} = \pi \circ f\) (where \(\pi : E \to E_{1}\) is the canonical homomorphism) belong to \(\widetilde{\mathcal{K}}(X;E_{1})\) , in which case \(\int f d\mu = \int f_{1} d\mu\) . We may therefore limit ourselves to considering only Hausdorff locally convex spaces.

\begin{enumerate}
\def\labelenumi{\arabic{enumi})}
\setcounter{enumi}{1}
\tightlist
\item
  Let E be a locally convex space over C, and let \(E_{0}\) be the locally convex space over R underlying E; we know that the mapping \(z^{\prime} \mapsto Rz^{\prime}\) which, to every continuous (complex) linear form \(z^{\prime}\) on E, makes correspond the continuous (real) linear form \(z \mapsto R\langle z, z^{\prime} \rangle\) on \(E_{0}\) , is an R-isomorphism of the dual \(E^{\prime}\) onto the dual \(E_{0}^{\prime}\) of \(E_{0}\) (TVS, II, §8, No.~1). Similarly, the algebraic dual \(E_{0}^{\prime *}\) of the real vector space \(E_{0}^{\prime}\) may be canonically identified with the real space underlying the algebraic dual \(E^{\prime *}\) of \(E^{\prime}\) . It follows that if \(\mu\) is a real measure and f a mapping in \(\widetilde{\mathcal{K}}(X;E)\) , the formula (1) is again valid when f is regarded as taking its values in \(E_{0}\) and the canonical bilinear forms figuring in the two members as being, respectively, relative to the duality between \(E_{0}^{\prime}\) and \(E_{0}^{\prime *}\) for the first member and the duality between \(E_{0}\) and \(E_{0}^{\prime}\) for the second.
\end{enumerate}

